% ============================================
% Section: Introduction
% ============================================
Photo-realistic 3D mapping of large or complex environments from a single
moving platform has matured rapidly. Systems such as Gaussian-LIC2
\cite{lang2026gaussian2} couple continuous-time LiDAR-Inertial-Visual (LIV)
odometry with 3D Gaussian Splatting (3DGS) \cite{kerbl20233dgs} to estimate an
accurate trajectory and construct a photo-realistic Gaussian map in real time.
A single platform, however, is fundamentally limited: it covers ground slowly,
cannot observe occluded regions from behind, and offers no redundancy against
sensor failure. Many realistic deployment scenarios---large-scale urban
surveying, multi-robot disaster response, autonomous warehouse fleets---require
\emph{multiple cooperative platforms} that share observations and maintain a
\emph{single} shared scene model.

Extending a single-platform LIV Gaussian SLAM system to a cooperative multi-node
setting raises three intertwined challenges that, to our knowledge, no prior
work addresses jointly:

\begin{enumerate}
  \item \textbf{Multi-sensor fusion at each node.} Real cooperative teams are
    heterogeneous: a survey vehicle may carry one LiDAR and one camera, while a
    drone may carry two cameras and a solid-state LiDAR. A unified fusion
    front-end must ingest an arbitrary number of cameras and LiDARs against a
    single continuous-time trajectory, with online estimation of every
    inter-sensor time offset.
  \item \textbf{Bandwidth-constrained data sharing.} A 3DGS map with
    $\sim$10\textsuperscript{6} Gaussians is far too large to stream raw over
    wireless links. Nodes must exchange compact, spatially bounded units and
    reconcile them into a consistent global map.
  \item \textbf{Global consistency under intermittent connectivity.} Links drop,
    nodes join and leave, and clocks drift. The system must produce a globally
    aligned photo-realistic map despite these disturbances, and degrade
    gracefully to single-node operation when peers disappear.
\end{enumerate}

\paragraph{Contributions.} We make the following contributions:

\begin{itemize}
  \item A \emph{generalized continuous-time multi-camera, multi-LiDAR, multi-IMU
    fusion front-end} (Section~\ref{sec:system}) that extends Coco-LIC
    \cite{lang2023coco} and Gaussian-LIC2 to $N_c$ cameras and $N_l$ LiDARs per
    node, with online time-offset estimation across all sensor pairs.
  \item A \emph{hierarchical hybrid communication architecture}
    (Section~\ref{sec:communication}) with cluster-head election, a submap
    publication protocol, and a bandwidth budget controller that keeps per-node
    uplink within a configurable cap while preserving map quality.
  \item A \emph{cooperative mapping back-end} (Section~\ref{sec:mapping}) that
    aligns incoming submaps via Gaussian-to-Gaussian ICP with descriptor and
    feature-based fallbacks, merges ownership-tagged Gaussians with confidence
    weighting, and maintains global consistency through a robust pose graph.
  \item A \emph{failure-mode analysis} (Section~\ref{sec:discussion}) covering
    head loss, link jitter, total partition, clock desync, and Byzantine nodes,
    with documented detection and recovery responses.
\end{itemize}

The full system design, including module APIs, configuration parameters, and
interface contracts, is documented in a companion design document
(\texttt{latex/design/DESIGN.md}). This report focuses on the system model,
protocol, and algorithms, and is intended as the basis for an experimental
evaluation campaign.

\paragraph{Scope.} This is a \emph{system design and algorithm report}. We
describe the architecture and analyze its expected behavior analytically and
through a bandwidth model; a full empirical evaluation on physical hardware is
left for future work (Section~\ref{sec:discussion}).
